\documentclass[border=10pt]{standalone}
\usepackage{tikz,ctex}
\usetikzlibrary{positioning}

\begin{document}

\begin{tikzpicture}[
    % 全局设置
    node distance = 1.6cm and 1.4cm,
    every node/.style = {
        draw,
        rounded corners,
        thick,
        align = center,
        font = \sffamily\small,
        minimum width = 2.0cm,
        minimum height = 0.9cm
    },
    % 样式定义
    centerStyle/.style = {
        fill = blue!45,
        text = white,
        font = \sffamily\bfseries\large,
        line width = 2pt,
        minimum size = 2.8cm,
        circle
    },
    branchTitle/.style = {
        font = \sffamily\bfseries,
        line width = 1.5pt
    },
    subNode/.style = {
        fill = white,
        font = \sffamily\small,
        minimum width = 2.2cm,
        minimum height = 0.8cm
    },
    leafNode/.style = {
        fill = white,
        font = \sffamily\footnotesize,
        minimum width = 1.9cm,
        minimum height = 0.7cm
    },
    arrowStyle/.style = {
        ->,
        thick,
        >=stealth,
        shorten >=2pt
    }
]

    % ================= 1. 中心节点 =================
    \node[centerStyle] (center) {混合密度网络 \\ Mixture Density \\ Network};

    % ================= 2. 上方分支：动机与逆问题 (红) =================
    \node[branchTitle, fill=red!15, above=of center] (motivation) {动机与逆问题 \\ Motivation \& Inverse Problems};

    \node[subNode, above=of motivation, xshift=-8cm] (gaussianlimit) {高斯假设局限 \\ Gaussian Limitation};
    \node[subNode, above=of motivation, xshift=-4cm] (robot) {机器人运动学 \\ Robot Kinematics};
    \node[subNode, above=of motivation, xshift=0cm] (forwardinverse) {正问题与逆问题 \\ Forward vs Inverse};
    \node[subNode, above=of motivation, xshift=4cm] (multimodal) {多模态分布 \\ Multimodal Distribution};
    \node[subNode, above=of motivation, xshift=8cm] (toyexample) {玩具逆问题 \\ Toy Inverse Problem};

    % ================= 3. 右侧分支：条件混合分布 (蓝) =================
    \node[branchTitle, fill=orange!15, right=of center] (mixture) {条件混合分布 \\ Conditional Mixture};

    \node[subNode, right=of mixture, yshift=4cm] (gmm) {高斯混合模型 \\ GMM (6.38)};
    \node[subNode, right=of mixture, yshift=2cm] (hetero) {异方差模型 \\ Heteroscedastic};
    \node[subNode, right=of mixture, yshift=0cm] (isotropic) {各向同性协方差 \\ Isotropic Covariance};
    \node[subNode, right=of mixture, yshift=-2cm] (cholesky) {Cholesky 分解扩展 \\ Cholesky Extension};
    \node[subNode, right=of mixture, yshift=-4cm] (experts) {与专家混合对比 \\ Mixture of Experts};

    % ================= 4. 下方分支：网络输出与误差函数 (绿) =================
    \node[branchTitle, fill=green!15, below=of center] (output) {网络输出与误差函数 \\ Outputs \& Error Function};

    \node[subNode, below=of output, xshift=-8cm] (picoef) {混合系数输出 \\ Softmax (6.40)};
    \node[subNode, below=of output, xshift=-4cm] (sigma) {方差输出 \\ Exp (6.41)};
    \node[subNode, below=of output, xshift=0cm] (mu) {均值输出 \\ Identity (6.42)};
    \node[subNode, below=of output, xshift=4cm] (nparams) {参数总数 \\ $(L+2)K$ Outputs};
    \node[subNode, below=of output, xshift=8cm] (neglog) {负对数似然误差 \\ NLL (6.43)};

    % ================= 5. 左侧分支：梯度优化与预测分布 (紫) =================
    \node[branchTitle, fill=purple!15, left=of center] (gradient) {梯度优化与预测分布 \\ Gradient \& Predictive Distribution};

    \node[subNode, left=of gradient, yshift=4cm] (gamma) {后验责任 \\ $\gamma_{nk}$ (6.44)};
    \node[subNode, left=of gradient, yshift=2cm] (gradpi) {混合系数梯度 \\ $\partial E_n/\partial a_k^\pi$ (6.45)};
    \node[subNode, left=of gradient, yshift=0cm] (gradmu) {均值梯度 \\ $\partial E_n/\partial a_{kl}^\mu$ (6.46)};
    \node[subNode, left=of gradient, yshift=-2cm] (gradsigma) {方差梯度 \\ $\partial E_n/\partial a_k^\sigma$ (6.47)};
    \node[subNode, left=of gradient, yshift=-4cm] (predictive) {预测分布 \\ Mean / Mode (6.48--6.50)};

    % ================= 连线逻辑 =================
    % 中心 -> 四大分支标题
    \draw[arrowStyle, red] (center) -- (motivation);
    \draw[arrowStyle, blue] (center) -- (mixture);
    \draw[arrowStyle, green!60!black] (center) -- (output);
    \draw[arrowStyle, purple] (center) -- (gradient);

    % 分支标题 -> 子节点 (上方)
    \draw[arrowStyle, red!60] (motivation) -- (gaussianlimit.south);
    \draw[arrowStyle, red!60] (motivation) -- (robot.south);
    \draw[arrowStyle, red!60] (motivation) -- (forwardinverse.south);
    \draw[arrowStyle, red!60] (motivation) -- (multimodal.south);
    \draw[arrowStyle, red!60] (motivation) -- (toyexample.south);

    % 分支标题 -> 子节点 (右侧)
    \draw[arrowStyle, blue!60] (mixture) -- (gmm.west);
    \draw[arrowStyle, blue!60] (mixture) -- (hetero.west);
    \draw[arrowStyle, blue!60] (mixture) -- (isotropic.west);
    \draw[arrowStyle, blue!60] (mixture) -- (cholesky.west);
    \draw[arrowStyle, blue!60] (mixture) -- (experts.west);

    % 分支标题 -> 子节点 (下方)
    \draw[arrowStyle, green!60!black] (output) -- (picoef.north);
    \draw[arrowStyle, green!60!black] (output) -- (sigma.north);
    \draw[arrowStyle, green!60!black] (output) -- (mu.north);
    \draw[arrowStyle, green!60!black] (output) -- (nparams.north);
    \draw[arrowStyle, green!60!black] (output) -- (neglog.north);

    % 分支标题 -> 子节点 (左侧)
    \draw[arrowStyle, purple!60] (gradient) -- (gamma.east);
    \draw[arrowStyle, purple!60] (gradient) -- (gradpi.east);
    \draw[arrowStyle, purple!60] (gradient) -- (gradmu.east);
    \draw[arrowStyle, purple!60] (gradient) -- (gradsigma.east);
    \draw[arrowStyle, purple!60] (gradient) -- (predictive.east);

\end{tikzpicture}

\end{document}